\documentclass[12pt]{article}
\usepackage{lmodern}
\usepackage{geometry}
\usepackage{enumitem}
\setlist[description]{leftmargin=1cm,labelindent=1cm}
\usepackage{amsmath}
\usepackage{amsthm}
\usepackage{amssymb}
\usepackage{mathtools}
\usepackage[portuguese]{babel}

\theoremstyle{plain}
\newtheorem{thm}{Teorema}[section]
\newtheorem{prop}[thm]{Proposição}
\newtheorem{lem}[thm]{Lema}
\newtheorem{cor}[thm]{Corolário}

\theoremstyle{definition}
\newtheorem{defn}[thm]{Definição}
\newtheorem{exmp}[thm]{Exemplo}
\newtheorem{prblm}[thm]{Problema}
\newtheorem{prcp}[thm]{Princípio}

\theoremstyle{remark}
\newtheorem*{rem}{Nota}

\geometry{
 a4paper,
 left=3cm,
 top=2cm,
 right=2cm,
 bottom=2cm
}

\begin{document}
\begin{center}
    \textbf{\large UNIVERSIDADE FEDERAL DO RIO GRANDE DO SUL} \\
    \textbf{\large INSTITUTO DE MATEMÁTICA E ESTATÍSTICA} \\
    \vspace{3cm}
    \textbf{\Large ÁLGEBRA LINEAR} \\
    \vspace{1cm}
    \textbf{\large LISTA DE RESULTADOS} \\
    \vspace{3cm}

    \textbf{LEONARDO SUAREZ}

    \vfill
    \textbf{\large Porto Alegre, RS} \\
    \textbf{\large \today}
\end{center}

\newpage
\tableofcontents
\newpage

\section{Espaços Vetoriais}
\section{Espaços Vetoriais de Dimensão Finita}

\begin{prop} (span é o menor subespaço que contém todos os vetores)
    O span de uma lista de vetores em \(V\) é o menor subespaço de \(V\) que contém todos os vetores da lista.
\end{prop}

\begin{defn}
    Um espaço vetorial \(V\) é dito de \emph{dimensão finita} se existe uma lista finita de vetores em \(V\) que gera \(V\). Caso contrário, \(V\) é dito de \emph{dimensão infinita}.
\end{defn}

\begin{defn}
    Uma lista \(v_1, \ldots, v_n\) de vetores em \(V\) é dita \emph{linearmente independente} se a única escolha de escalares \(a_1, \ldots, a_n\) em \(\mathbb{F}\) que satisfaz
    \begin{equation}
        a_1 v_1 + a_2 v_2 + \cdots + a_n v_n = 0
    \end{equation}
    \noindent é \(a_1 = a_2 = \cdots = a_n = 0\). Caso contrário, a lista é dita \emph{linearmente dependente}. A lista vazia é dita linearmente independente.
\end{defn}

\begin{lem} (lema de dependência linear)
    Suponha que \(v_1, \ldots, v_n\) é linearmente dependente em \(V\). Então existe \(j \in \{1, \ldots, n\}\) tal que
    \begin{equation}
        v_j \in \text{span}(v_1, \ldots, v_{j-1}).
    \end{equation}
    \noindent Ainda, se \(j\) satisfaz a condição acima e o \(j\)-ésimo vetor for removido da lista, então o span dos vetores restantes é igual a \(\text{span}(v_1, \ldots, v_n)\).
\end{lem}

\begin{prop} (tamanho de uma lista linearmente independente \(\leq\) tamanho de uma lista geradora)
    Em um espaço vetorial \(V\) de dimensão finita, o tamanho de qualquer lista linearmente independente é menor ou igual ao tamanho de qualquer lista geradora.
\end{prop}

\begin{prop} (subespaços de dimensão finita)
    Todo subespaço de um espaço vetorial de dimensão finita é de dimensão finita.
\end{prop}

\begin{defn}
    Uma \emph{base} de um espaço vetorial \(V\) é uma lista de vetores \(v_1, \ldots, v_n\) em \(V\) que é linearmente independente e gera \(V\).
\end{defn}

\begin{prop} (critério para base)
    Uma lista de vetores \(v_1, \ldots, v_n\) em um \(V\) é uma base de \(V\) sse todo vetor em \(V\) pode ser escrito de maneira única como uma combinação linear de \(v_1, \ldots, v_n\), isto é,
    \begin{equation}
        v = a_1 v_1 + a_2 v_2 + \cdots + a_n v_n,
    \end{equation}
    \noindent onde \(a_1, \ldots, a_n\) são escalares em \(\mathbb{F}\).
\end{prop}

\begin{prop} (toda lista geradora possui uma base)
    Toda lista geradora em um espaço vetorial pode ser reduzida a uma base do espaço vetorial.
\end{prop}

\begin{prop} (toda lista linearmente independente pode ser estendida a uma base)
    Toda lista linearmente independente em um espaço vetorial de dimensão finita pode ser estendida a uma base do espaço vetorial.
\end{prop}

\begin{prop} (todo subespaço de \(V\) é parte de uma soma direta igual a \(V\))
    Suponha que \(V\) seja um espaço vetorial de dimensão finita e \(U\) é um subespaço de \(V\). Então existe um subespaço \(W\) de \(V\) tal que \(V = U \oplus W\).
\end{prop}

\begin{prop} (o tamanho da base não depende da base)
    Se \(V\) é um espaço vetorial de dimensão finita, então todas as bases de \(V\) têm o mesmo tamanho.
\end{prop}

\begin{defn}
    A \emph{dimensão} de um espaço vetorial \(V\) de dimensão finita é o tamanho de qualquer base de \(V\).
\end{defn}

\begin{prop} (dimensão do subespaço)
    Se \(V\) é um espaço vetorial de dimensão finita e \(U\) é um subespaço de \(V\), então \(\dim(U) \leq \dim(V)\).
\end{prop}

\begin{prop} (lista linearmente independente com tamanho igual à dimensão é base)
    Se \(V\) é um espaço vetorial de dimensão finita, então toda lista linearmente independente de vetores em \(V\) com tamanho igual a \(\dim(V)\) é uma base de \(V\).
\end{prop}

\begin{cor} (subespaço com mesma dimensão é o espaço todo)
    Se \(V\) é um espaço vetorial de dimensão finita, então todo subespaço \(U\) de \(V\) com \(\dim(U) = \dim(V)\) é igual a \(V\).
\end{cor}

\begin{prop} (lista geradora com tamanho igual à dimensão é base)
    Se \(V\) é um espaço vetorial de dimensão finita, então toda lista geradora de vetores em \(V\) com tamanho igual a \(\dim(V)\) é uma base de \(V\).
\end{prop}

\begin{prop} (dimensão de soma)
    Se \(V_1\) e \(V_2\) são subespaços de um espaço vetorial de dimensão finita \(V\), então
    \begin{equation}
        \dim(V_1 + V_2) = \dim(V_1) + \dim(V_2) - \dim(V_1 \cap V_2).
    \end{equation}
\end{prop}

\section{Transformações Lineares}

\begin{defn}
    Uma transformação linear de \(V\) em \(W\) é uma função \(T: V \to W\) com as seguintes propriedades.
    \begin{enumerate}
        \item \(T(u + v) = T(u) + T(v)\), para todo \(u, v \in V\).\hfill\emph{(aditividade)}
        \item \(T(\lambda u) = \lambda T(u)\), para todo \(u \in V\) e todo escalar \(\lambda \in \mathbb{F}\).\hfill\emph{(homogeneidade)}
    \end{enumerate}
\end{defn}

\begin{lem}
    Suponha que \(v_1,\ldots,v_n\) é uma base de \(V\) e \(w_1,\ldots,w_n \in W\). Então existe uma única transformação linear \(T: V \to W\) tal que \(T(v_k) = w_k\), para \(k = 1, \ldots, n\).
\end{lem}

\begin{defn}
    Suponha \(S, T \in \mathcal{L}(V, W)\) e \(\lambda \in \mathbb{F}\). A \emph{soma} \(S + T\) e o \emph{produto por escalar} \(\lambda T\) são as transformações lineares de \(V\) em \(W\) definidas por
    \[
        (S + T)(v) = S(v) + T(v) \quad \text{e} \quad (\lambda T)(v) = \lambda (T(v)),
    \]
    \noindent para todo \(v \in V\).
\end{defn}

\begin{cor}
    Com as operações de soma e produto escalar definidas acima, \(\mathcal{L}(V, W)\) é um espaço vetorial sobre \(\mathbb{F}\).
\end{cor}

\begin{defn}
    Se \(T \in \mathcal{L}(U, V)\) e \(S \in \mathcal{L}(V, W)\), então o \emph{produto} \(ST\) é a transformação linear de \(U\) em \(W\) definida por
    \[
        (ST)(u) = S(T(u)),
    \]
    \noindent para todo \(u \in U\).
\end{defn}

\begin{prop} (propriedades algébricas do produto de transformações lineares)
    \begin{itemize}
        \item \(\left(T_1 T_2\right)T_3 = T_1\left(T_2T_3\right)\) sempre que \(T_1, T_2, T_3\) são transformações lineares tais que os produtos são definidos.\hfill\emph{(associatividade)}
        \item \(T_1I = IT_1 = T_1\) sempre que \(T_1 \in \mathcal{L}(V, W)\), onde o primeiro \(I\) é a identidade em \(V\) e o segundo \(I\) é a identidade em \(W\).\hfill\emph{(elemento neutro)}
        \item \(\left(S_1 + S_2\right)T = S_1T + S_2T\) e \(S\left(T_1 + T_2\right) = ST_1 + ST_2\), sempre que \(T, T_1, T_2 \in \mathcal{L}(U, V)\) e \(S, S_1, S_2 \in \mathcal{L}(V, W)\).\hfill\emph{(distributividade)}
    \end{itemize}
\end{prop}

\begin{cor}
    Suponha que \(T\) seja uma transformação linear de \(V\) em \(W\). Então \(T(0) = 0\).
\end{cor}

\begin{defn}
    Para \(T \in \mathcal{L}(V, W)\), o \emph{espaço nulo} de \(T\) é o subconjunto de \(V\) definido por
    \[
        \mathcal{N}(T) = \{v \in V : T(v) = 0\}.
    \]
\end{defn}

\begin{prop}
    Seja \(T \in \mathcal{L}(V, W)\). Então \(\mathcal{N}(T)\) é um subespaço de \(V\).
\end{prop}

\begin{defn}
    Uma função \(T: V \to W\) é \emph{injetora} se \(T(v_1) = T(v_2)\) implica \(v_1 = v_2\), para todo \(v_1, v_2 \in V\).
\end{defn}

\begin{prop}
    Seja \(T \in \mathcal{L}(V, W)\). Então \(T\) é injetora se e somente se \(\mathcal{N}(T) = \{0\}\).
\end{prop}

\begin{defn}
    Para \(T \in \mathcal{L}(V, W)\), a \emph{imagem} de \(T\) é o subconjunto de \(W\) definido por
    \[
        \text{Im}(T) = \{T(v) : v \in V\}.
    \]
\end{defn}

\begin{prop}
    Seja \(T \in \mathcal{L}(V, W)\). Então \(\text{Im}(T)\) é um subespaço de \(W\).
\end{prop}

\begin{defn}
    Uma função \(T: V \to W\) é \emph{sobrejetora} se para todo \(w \in W\) existe \(v \in V\) tal que \(T(v) = w\).
\end{defn}

\begin{thm} (teorema fundamental das transformações lineares)
    Suponha que \(V\) seja um espaço vetorial de dimensão finita e que \(T \in \mathcal{L}(V, W)\). Então \(\text{Im}(T)\) é de dimensão finita e
    \begin{equation}
        \dim V = \dim \mathcal{N}(T) + \dim \text{Im}(T).
    \end{equation}
\end{thm}

\begin{cor}
    Suponha que \(V\) e \(W\) são espaços vetoriais de dimensão finita com \(\dim V > \dim W\). Então nenhuma transformação linear de \(V\) em \(W\) é injetora.
\end{cor}

\begin{cor}
    Suponha que \(V\) e \(W\) são espaços vetoriais de dimensão finita com \(\dim V < \dim W\). Então nenhuma transformação linear de \(V\) em \(W\) é sobrejetora.
\end{cor}

\section{Polinômios}
\section{Autovalores e Autovetores}
\section{Espaços com Produto Interno}
\section{Operadores em Espaços com Produto Interno}
\section{Operadores em Espaços Vetoriais Complexos}
\section{Álgebra Multilinear e Determinantes}

\end{document}